\documentclass[11pt]{article}
\usepackage[margin=2.4cm]{geometry}
\usepackage{graphicx}
\usepackage{booktabs}
\usepackage{amsmath}
\usepackage{microtype}
\usepackage[round]{natbib}
\usepackage[colorlinks=true,linkcolor=black,citecolor=black,urlcolor=blue]{hyperref}
\usepackage{caption}
\captionsetup{font=small,labelfont=bf}
\setlength{\parskip}{2pt}

\title{Two-frame weighted PCA precleaning of near-infrared spectra,\\
and what it does to line-by-line velocities}
\author{É. Artigau}
\date{Draft of 2026 September 12}

\begin{document}
\maketitle

\begin{abstract}
Near-infrared radial velocities are limited less by photon noise than by what
is left of the atmosphere and of the instrument in a corrected spectrum. We
model a campaign with two rest frames at once: one anchored on the star, one
anchored on the observer, fitted together on a log-uniform grid where a Doppler
shift is an exact translation. Dividing the observer block out of every
exposure leaves a spectrum whose stellar content is untouched, which is then
measured line by line. On a SPIRou campaign of TOI-2120, where telluric
residuals dominate, the velocity scatter falls from 47.7 to 17.4~m\,s$^{-1}$.
On NIRPS campaigns whose delivered velocities are already at 2.7~m\,s$^{-1}$ the
same correction moves them by less than a tenth of a m\,s$^{-1}$, in either
direction. A pair of twin experiments, two 14-night slices
of one campaign at the same signal-to-noise, shows what decides the outcome:
where the barycentric velocity spans 42.8~km\,s$^{-1}$ the correction improves
the robust scatter from 8.5 to 4.5~m\,s$^{-1}$, and where it spans
0.7~km\,s$^{-1}$ it degrades it from 9.1 to 21.8. When the star does not move
against the sky, the two frames are one frame, the observer block takes up
stellar structure, and dividing it out removes part of the star. We give a
practical guard, and a mode that fits several stars against a single observer
basis, which breaks the degeneracy by construction.
\end{abstract}

\section{Introduction}

In the near infrared the error budget of a precise radial velocity is rarely
photon noise. Telluric absorption and emission, detector systematics and the
imperfect removal of both leave a residual that is fixed in the observer's
frame while the star moves through it at the barycentric velocity. Telluric
correction removes most of it; what remains is small in flux and large in
velocity.

The line-by-line method of \citet{artigau2022} measures a velocity per line and
combines them with an outlier-resistant estimator, which is what makes a
residual of this kind visible: a handful of contaminated lines no longer drags
the whole exposure. It also makes it worth removing the residual before the
velocities are measured rather than after.

This paper describes a data-driven precleaning of that residual, and, as much
as the method itself, what four campaigns say about when it helps. The
implementation, the configuration of every run shown here and the tests that
pin its behaviour are in one repository; each figure and each number below is
produced by the pipeline it documents.

\section{Method}

\subsection{Two frames, one model}

Each exposure is resampled onto a grid uniform in $\ln \lambda$, where a
Doppler shift is a translation of $\mathrm{atanh}(v/c)$ divided by the grid's
step, exactly, with no approximation in $v/c$. A Savitzky-Golay filter of
100~km\,s$^{-1}$ removes the continuum, so the data are
$y = \ln f - \mathrm{savgol}(\ln f)$, in which an absorption of optical depth
$\tau$ is $-\tau$ and scales linearly with the column it comes from.

A t.fits exposure gives two rows, the even orders and the odd ones, which see
the same lines at different resolutions. The model of row $n$ is
%
\begin{equation}
  y_n \;=\; T_{p(n)} \;+\; a_n \, S_n P \;+\; b_n \cdot Q \;+\; \epsilon_n ,
  \label{eq:model}
\end{equation}
%
where $T_p$ is the star's own spectrum for order parity $p$, carried into the
exposure's frame by $S_n$, the shift operator of that exposure's barycentric
velocity; $P$ are components anchored in the star's rest frame; $Q$ are
components anchored in the observer's; and $a_n$, $b_n$ are that exposure's
amplitudes, shared by its two parity rows because they are one measurement of
one star at one instant.

Two properties of Eq.~\ref{eq:model} matter for what follows.

\paragraph{The star's mean spectrum carries the coefficient one.} $T_p$ is
estimated once, before any component, as the barycentric-binned median of the
rows of that parity in the star's frame, and is then subtracted with a
coefficient of exactly one. A coefficient in front of a term in a logarithm is
an exponent on the flux, and an exponent on a star's mean spectrum describes no
star. The free amplitude $a_n$ multiplies $P$, which describes variability
around that fixed spectrum, and is measured below to be nearly orthogonal to
it.

\paragraph{A free amplitude on an observer component is exact physics.} For an
absorption, $\ln T = -\tau$ and the optical depth scales with the airmass and
the column, so $b_n$ is the physical parameter itself, not a linearisation.

\subsection{The star per order parity}

The two halves of the echellogram are not equivalent. Where two orders overlap,
the same stellar lines appear at both ends of a blaze, and their measured
depths differ: on Proxima at 1201.1~nm the four lines of one overlap are
shallower in the odd orders by 0.005 to 0.042 in $\ln f$, and the even and odd
residuals of a single fit are mirror images of one another, with a rank
correlation of $-0.98$. One star spectrum cannot describe both halves. The
model therefore carries one $T_p$ per parity, and only the amplitudes are
shared.

\subsection{Weights, and the samples nobody should trust}

Every sample carries a weight: photon noise, floored by the sample-to-sample
scatter measured over 201 samples, which is what keeps the thermal background
of the reddest orders from being called the quietest part of the detector; a
ramp in telluric transmission; zero below a line depth of $\ln f = -0.5$; zero
at the edges of the grid; and zero wherever the Lanczos taps of the carry would
reach into a hole in the basis support.

\subsection{Solving}

The two blocks are solved by block coordinate descent. Each sweep solves the
amplitudes of both blocks jointly per exposure, by the normal equations of the
design $[\,S_n P^{\mathsf T} \mid Q^{\mathsf T}\,]$, the two parity rows
summed into one system; then updates each basis against the residual of the
other. The star side is a cubic B-spline with a knot per sample, evaluated at
each exposure's shifted positions and updated exactly by a banded solve, which
is both faster and slightly better conditioned than carrying samples with a
Lanczos kernel.

\subsection{The correction}

Only the observer block is divided out:
%
\begin{equation}
  f_{\rm corrected} \;=\; f \, \exp\!\left(- b_n \cdot \tilde{Q} \right) ,
\end{equation}
%
on each order's own pixels, by a cubic spline from the grid. The star block is
never divided out: with no external template the first star component is the
star.

$\tilde{Q}$ is the fitted basis shrunk by its own significance, column by
column. At column $i$ the components' Fisher matrix is
$F_i = \sum_n w_{ni} b_n b_n^{\mathsf T}$, with the weights scaled by the fit's
median reduced $\chi^2$; $\sigma^2(Q_{ji}) = [F_i^{-1}]_{jj}$ and
$z_{ji} = Q_{ji}/\sigma(Q_{ji})$; and
%
\begin{equation}
  \tilde{Q}_{ji} \;=\; \max\!\left(0,\, 1 - z_{ji}^{-2}\right) Q_{ji} .
\end{equation}
%
The significance is the basis's, measured with every exposure at once, not one
exposure's share of it: a pattern at half a sigma in each of $N$ exposures is
detected at about $0.5\sqrt{N}$ and is kept, while a column where the component
is noise is divided out at zero and the file keeps its delivered flux there.

Finally, a sample the fit gave no weight cannot be corrected, and is blanked.
Blanking each exposure's own leaves every epoch with a different set of lines;
we show below that this alone costs 1.8~m\,s$^{-1}$ on one campaign. The
nominal mode blanks, in every exposure, every sample that any exposure lost, so
the campaign carries one set of lines.

\section{Data}

Four campaigns, two instruments, reduced with APERO \citep{cook2022} and
measured with LBL \citep{artigau2022}.

\begin{table}[h]
\centering
\small
\begin{tabular}{lllrrrrr}
\toprule
Target & Instrument & $T_{\rm eff}$ (K) & exposures & nights & $t_{\rm exp}$ (s)
& SNR$_H$ & delivered rms (m\,s$^{-1}$)\\
\midrule
TOI-2120 & SPIRou & 3179 & 321 & --- & 903 &  51 & 47.7\\
Proxima  & NIRPS  & 3050 & 782 & 259 & 201 & 262 &  2.9\\
GJ~1     & NIRPS  & 3478 & 292 & 146 & 178 & 217 &  2.7\\
TOI-4552 & NIRPS  & 3304 & 119 & 60  & 897 &  43 & 12.6\\
\bottomrule
\end{tabular}
\caption{The campaigns. Nights are the nightly coadds the NIRPS fits use;
$t_{\rm exp}$ and SNR$_H$ are medians over every exposure of the campaign, the
latter of the per-order extraction signal-to-noise of the orders whose centres
fall in 1490--1780~nm, a band both instruments cover, so that the two can be
compared; the delivered rms is the line-by-line velocity scatter of the spectra
as reduced, with no precleaning.}
\label{tab:campaigns}
\end{table}

\section{Results}

\subsection{What the correction is worth}

\begin{table}[h]
\centering
\small
\begin{tabular}{lrrrrr}
\toprule
& \multicolumn{2}{c}{delivered} & \multicolumn{2}{c}{corrected} & \\
\cmidrule(lr){2-3}\cmidrule(lr){4-5}
Target & rms & robust $\sigma$ & rms & robust $\sigma$ & exposures\\
\midrule
TOI-2120 & 47.71 & 34.99 & 17.40 & 14.07 & 316\\
Proxima  &  2.86 &  2.22 &  2.93 &  2.37 & 765\\
GJ~1     &  2.70 &  2.74 &  2.61 &  2.38 & 280\\
TOI-4552 & 12.64 & 11.37 & 14.24 & 11.20 & 109\\
\bottomrule
\end{tabular}
\caption{Line-by-line velocities, in m\,s$^{-1}$, on the exposures both series
share, with the adopted parameters: no star component, three observer
components, the correction shrunk by its significance, one set of lines per
campaign, and each series measured against the template LBL builds from it. The
correction transforms the campaign whose residuals are telluric. Where the
delivered velocities are already at 2.7~m\,s$^{-1}$ it neither helps nor costs
much, and on GJ~1 it takes the robust scatter below the delivered series.}
\label{tab:overall}
\end{table}

The pattern of Table~\ref{tab:overall} is not a signal-to-noise effect, and the
cleanest case against that reading is TOI-2120 against TOI-4552. Their exposures
are of the same length, 903 and 897~s, and their spectra reach the same
signal-to-noise in the band both instruments cover, 51 and 43 per order in $H$
(Table~\ref{tab:campaigns}), yet the correction removes nearly two thirds of the
velocity scatter of one and adds a sixth to the other's rms while leaving its
robust scatter where it was. The two brighter campaigns, whose spectra have four
to six times that signal-to-noise, are the ones where the correction has the
least to remove, and there it neither helps nor costs more than a tenth of a
m\,s$^{-1}$.

\subsection{The barycentric span decides}

TOI-4552 was observed in two seasons. In the first the barycentric velocity
spans 52.6~km\,s$^{-1}$; in the second, 9.8. Fitting each season alone, the
correction improves the first and degrades the second. To remove every other
difference we then cut two slices of 14 nights each, at the same median
signal-to-noise, one spanning 42.8~km\,s$^{-1}$ of barycentric velocity and one
spanning 0.7 (Fig.~\ref{fig:berv}):

\begin{table}[h]
\centering
\small
\begin{tabular}{lrrrr}
\toprule
& \multicolumn{2}{c}{delivered} & \multicolumn{2}{c}{corrected}\\
\cmidrule(lr){2-3}\cmidrule(lr){4-5}
14 nights of TOI-4552 & rms & robust $\sigma$ & rms & robust $\sigma$\\
\midrule
BERV span 42.8 km\,s$^{-1}$ & 7.80 & 8.48 & 6.18 & 4.52\\
BERV span \phantom{0}0.7 km\,s$^{-1}$ & 8.96 & 9.14 & 21.76 & 21.76\\
\bottomrule
\end{tabular}
\caption{The same campaign, the same instrument, the same number of nights and
the same signal-to-noise. Only the barycentric coverage differs, by a factor
of sixty, and the correction changes sign.}
\end{table}

\begin{figure}[h]
\centering
\includegraphics[width=0.49\textwidth]{../outputs/TOI4552/rv_bervwide.pdf}
\includegraphics[width=0.49\textwidth]{../outputs/TOI4552/rv_bervnarrow.pdf}
\caption{Left: the wide-span slice, corrected in orange over the delivered
series in grey. Right: the narrow-span slice. Nightly means above, exposures
below.}
\label{fig:berv}
\end{figure}

\subsection{What does not explain it}

Every other candidate was measured on the same 109 exposures of TOI-4552 and
ruled out: nightly stacking (fitting the exposures individually gives 15.6
against 15.9~m\,s$^{-1}$); the number of observer components, in the fit
(16.9 with one instead of three) or in the correction (15.9 with one divided
out); the velocity term (16.4); the shrinkage, which helps (15.1 with it, 18.1
without); the template handed to LBL; and the mask.

\subsection{The mask, and the set of lines}

Blanking the samples the fit could not weight is necessary, but blanking a
different set in every exposure is not. With nothing divided out at all and
each exposure blanking its own, the scatter rises from 12.64 to
14.41~m\,s$^{-1}$ while the per-exposure error does not move: a line set that
moves from epoch to epoch is scatter, not lost information. Blanking one common
set instead gives 11.96, below the delivered spectra themselves, for 4.8
percentage points more of each exposure blanked.

\subsection{The star's freedom, and what $a_1$ is}

The amplitude $a_1$ multiplies the star-frame variability component, never the
star's mean spectrum. Measured on four campaigns, that component is nearly
orthogonal to the mean spectrum, $|\cos| \le 0.065$; expressed as an exponent
on the star it amounts to $0.4\%$ rms on the three NIRPS targets and $1.25\%$
on TOI-2120.

It is nevertheless not describing the star. On Proxima, $a_1$ follows the
barycentric velocity with a rank correlation of 0.69, ahead of humidity (0.43)
and airmass ($-0.39$), and its periodogram peaks at 366 days
(Fig.~\ref{fig:a1}). It is an annual mode aligned with the Earth's orbit.

Removing the star component altogether, so that the star is its fixed template
and nothing else, takes TOI-4552 from 14.56 to 11.20~m\,s$^{-1}$ of robust
scatter at equal mask and template (Fig.~\ref{fig:freedom}), that is, below the
delivered spectra. The freedom the star block is given is used, in part, to
absorb what belongs to the observer, and the observer block then divides out
what belongs to the star.

\begin{figure}[h]
\centering
\includegraphics[width=0.62\textwidth]{../outputs/PROXIMA/a1_series_proxima.pdf}
\caption{Proxima: the star amplitude over time, against the barycentric
velocity, and its periodogram. The right-hand axis of the first panel gives
the equivalent exponent on the star's mean spectrum.}
\label{fig:a1}
\end{figure}

\begin{figure}[h]
\centering
\includegraphics[width=0.92\textwidth]{../outputs/TOI4552/rv_star_freedom.pdf}
\caption{TOI-4552 at equal mask and template: the delivered series, the same
spectra with the mask alone, the correction with a free star amplitude, and
the correction with no star component at all.}
\label{fig:freedom}
\end{figure}

\section{Discussion}

\subsection{One frame instead of two}

The observer block is identifiable only because the star moves through it. Over
a campaign whose barycentric velocity spans tens of km\,s$^{-1}$, a pattern
fixed in the observer's frame lands on different stellar lines in every
exposure, and no combination of stellar structure can imitate it. Over a
campaign whose barycentric velocity barely moves, the two frames coincide: a
column of the observer basis and a column of the star's spectrum are then the
same column, and the fit has no way to tell which of them owns the flux.
Dividing the observer block out then removes part of the star, and the
line-by-line velocities record it as a barycentric-shaped injection.

This is the same mechanism, at the campaign scale, as the annual mode that
$a_1$ carries.

\subsection{A guard, and a way out}

The practical guard is to look at the barycentric span of the set of exposures
a fit is built on, not merely at their number, and to fit seasons separately
when a campaign mixes wide and narrow spans: on TOI-4552, fitting the narrow
season alone costs 14\% against the delivered velocities, where fitting it
together with the wide one costs 55\%.

The principled way out is to break the degeneracy with other stars. The
atmosphere and the instrument belong to the night, not to the target, while the
star, its barycentric coverage and its systemic velocity do not: several
campaigns fitted against a single observer basis cannot let that basis follow
any one star. The pipeline has a mode for it, in which the rows of several
objects share one grid and one observer basis while each object keeps its own
star spectrum per order parity.

\subsection{What the precleaning is for}

Table~\ref{tab:overall} is the honest summary: this is a method for campaigns
whose velocities are limited by what the atmosphere and the instrument left
behind, and it is not free where they are not. On TOI-2120 it removes nearly two
thirds of the scatter, 47.7 to 17.4~m\,s$^{-1}$. On Proxima and GJ~1 the
delivered velocities are already at 2.7 to 2.9~m\,s$^{-1}$ and the correction
moves them by less than a tenth of a m\,s$^{-1}$, upward on one and downward on
the other. Applying it there buys nothing and costs nothing, and the decision can
be made in advance from the barycentric span and from how much of the scatter the
formal errors already explain.

\section{Conclusions}

\begin{enumerate}
\item A two-frame weighted decomposition separates what is fixed in the
observer's frame from what belongs to the star, and dividing the first out
removes nearly two thirds of the velocity scatter of a campaign dominated by
telluric residuals, 47.7 to 17.4~m\,s$^{-1}$ on TOI-2120.
\item What decides whether it helps is the barycentric span of the fitted set.
Two 14-night slices of one campaign, at equal signal-to-noise, differ by a
factor of sixty in span, and the correction improves the robust scatter by a
factor of two in one and degrades it by a factor of 2.4 in the other.
\item The star's mean spectrum must carry a coefficient of exactly one, and the
star is better given no free component at all. The freedom given to one is used
to absorb an annual, barycentric mode; removing it improved every campaign
measured, by 0.9~m\,s$^{-1}$ on TOI-2120, 1.6 on TOI-4552 and 0.4 on GJ~1, which
it takes below its delivered velocities.
\item Blanking the samples a fit could not weight must be done with one set of
lines for the whole campaign. Blanking each exposure's own costs
1.8~m\,s$^{-1}$ with no correction applied at all.
\end{enumerate}

\begin{thebibliography}{9}
\bibitem[Artigau et al.(2022)]{artigau2022} Artigau, É., et al. 2022,
  \textit{Line-by-line Velocity Measurements: an Outlier-resistant Method for
  Precision Velocimetry}, AJ, 164, 84, doi:10.3847/1538-3881/ac7ce6
\bibitem[Cook et al.(2022)]{cook2022} Cook, N. J., et al. 2022,
  \textit{APERO: A PipelinE to Reduce Observations, Demonstration with SPIRou},
  PASP, 134, 114509, doi:10.1088/1538-3873/ac9e74
\end{thebibliography}

\end{document}
